\histitem{Locally compact commutative groups}

The interaction between harmonic analysis and group representations was thus firmly established by the early 1930s. It was greatly strengthened in the following decade, with each of the two subjects prompting rapid progress in the other.

A very fertile ground for this interaction was the study of almost periodic functions, introduced by H. Bohr in 1924 in the case of functions of a real variable {[}9{]}. By one of the fundamental theorems of the latter, these are the bounded continuous functions that are uniform limits on \(\mathbf{R}\) of linear combinations of functions \(x \mapsto e^{i\lambda x}\) , \(\lambda \in \mathbf{R}\) (cf. II, p. 292, exerc. 54). Bochner {[}6{]} had proved in 1926 that a function \(f \in \mathcal{C}_b(\mathbf{R})\) is almost periodic if and only if its translates \(x \mapsto f(x - a)\) , for \(a \in \mathbf{R}\) , form a relatively compact subset of \(\mathcal{C}_b(\mathbf{R})\) for the topology of uniform convergence. Bochner and von Neumann observed that if G is a topological group, this characterization visibly provides a notion of almost periodic function (on the right or on the left) on G. In the early 1930s, the theory promised to develop rapidly, in particular under the impetus of von Neumann {[}52{]} and the methods of Peter and Weyl.

Another fertile area appears on the periphery of N. Wiener's work. He introduced, in a famous paper {[}88{]} from 1932, algebraic methods in the study of so-called Tauberian problems—that is, in the study of the asymptotic behavior of a function (or of a sequence) along a filter, given information concerning the behavior of certain weighted means. His approach was motivated by an idea of R. Schmidt.

One of the auxiliary results of Wiener's work ({[}88, Lemma IIe{]}) will particularly strike the minds and inspire numerous developments: if \(f\) is a periodic function whose Fourier series is absolutely convergent, and if \(f\) does not vanish, then the Fourier series of \(1/f\) is absolutely convergent (cf. I, p. 38, example). As early as 1932, R. Paley and N. Wiener {[}89{]} sketched the link between this convergence result, the notion of dual group for a discrete commutative group, and the theory of almost periodic functions.

Against this analytic backdrop, it is algebraic topology that leads L. Pontryagin to give the decisive impetus for the study of characters of locally compact abelian groups. He announces as early as 1932 {[}55{]} the interest of the notion of the character group lies in placing in duality the homology groups of a compact subset of Euclidean space and those of its complement. In a paper {[}58{]} published in 1934, he shows that the character group of a discrete group is compact, and then establishes the duality theorems in this framework. Although Pontryagin uses the results of Peter and Weyl, as well as the Haar measure whose existence had just been established in general, he aims mainly at applications to topology {[}57{]} and is not explicitly concerned with harmonic analysis \(^{(8)}\): the duality theorems are proved by a fine study of the structure of compact abelian groups.

Pontryagin's work immediately aroused great interest. On a suggestion of von Neumann, E. van Kampen extended the duality the following year to all locally compact commutative groups {[}37{]}. Shortly thereafter, he gives a spectacular application of the new invariant analysis to the study of almost periodic functions {[}38{]}, discovered independently by A. Weil during the same period {[}79{]}: If G is a locally compact abelian group, one obtains a compact group \(\widetilde{G}\) beginning by endowing the dual group \(\widehat{G}\) with the discrete topology, and then defining \(\widetilde{G}\) as the compact dual group of the discrete group thus obtained. The group G is then canonically embedded in \(\widetilde{G}\) (cf. II, p. 292, exerc. 54), and a simple application of the Peter--Weyl theory for continuous functions on \(\widetilde{G}\) yields all the results known at the time about almost periodic functions on G.

At the same time, functions of positive type, well known in classical analysis, appear at this period in the study of unitary representations of R. The notion of a universally positive kernel had been studied by J. Mercer {[}45{]} shortly before 1910, in connection with the theory of integral equations, while its analogue for sequences gave rise to a considerable number of works during the same period, related in particular to the moment problem (cf. IV, p. 359, exerc. 21). Meanwhile, functions of positive type on R, already studied systematically by M. Mathias {[}43{]} in 1923, had become, in Bochner's hands, one of the fundamental tools of Fourier analysis and of probability theory {[}7{]}. In 1933, Bochner {[}8{]} and Riesz {[}61{]} independently observed that every diagonal matrix coefficient of a unitary representation of R is a function of positive type on R. This observation enabled them to give a simpler proof of the Stone theorem, and was to have a decisive influence on the development of the general theory (cf. n° 7).

The results of this period were organized by A. Weil in a book completed before 1937 and published in 1940 {[}80{]}. There he sets out in detail the results of Peter--Weyl, Pontryagin, and van Kampen. But in the case of commutative groups, whereas Pontryagin and van Kampen emphasized the structure of groups and duality theorems, Weil systematically develops harmonic analysis, introducing the Fourier transform and proving the Plancherel formula and Bochner's theorem (V, p. 455, th. 5). He insists in particular on the role of convolution (under the name "composition product") and on that of functions of positive type, now defined on an arbitrary group and related to the diagonal matrix coefficients of representations, at least in the case of finite-dimensional representations.

The generality thus achieved in the commutative case would open, a little later, new horizons in number theory. As early as 1936, with a view to introducing algebraic methods into class field theory, C. Chevalley {[}14{]} introduced the group of idèles of a number field, which would later appear as the group of invertible elements of the ring of adèles (cf. AC, VII, p. 221--222, and ÉHM, p. 143). In 1950, J. Tate {[}75{]} would show that harmonic analysis for groups of adèles and idèles makes it possible to recover easily the functional equations of Hecke L-functions, foreshadowing extraordinary developments linking harmonic analysis on adelic groups to the study of automorphic forms.

But one can say that from the end of the 1930s, the principal results of invariant harmonic analysis had been established for compact groups and for locally compact commutative groups. In the commutative case, Weil's proofs still depend (like those of Pontryagin and van Kampen) on a fine knowledge of the structure of the group, that is, roughly speaking, on the classification results of § 2 of II, p. 244. The following decade would show how to dispense with this, thanks to the new methods of the Gelfand school (cf. H. Cartan and R. Godement {[}12{]}).

